\section{Related Work}
\label{sec:related_work}
% RLHF is now widely regarded as an effective approach for aligning LLMs with human values and preferences \cite{ouyang2022training,casper2023open}. The conventional RLHF pipeline typically consists of three core stages: supervised fine-tuning, reward model construction and training, and policy optimization \cite{achiam2023gpt,touvron2023llama,liu2024deepseek}. In the policy optimization stage, commonly used algorithms include Proximal Policy Optimization (PPO) \cite{schulman2017proximal}, Group Relative Policy Optimization \cite{liu2024deepseek}, and REINFORCE++ \cite{hu2025reinforce++}, among others. Although RLHF has demonstrated impressive performance, it also complicates the training process by requiring multiple models and sampling from LLMs \cite{zhao2023slic, casper2023open,xu2024contrastive,chen2024bootstrapping}. Consequently, researchers have been actively exploring more efficient and streamlined offline alternatives \cite{rafailov2023direct,yuan2023rrhf,zhao2023slic,xu2024contrastive,ethayarajh2024kto,meng2024simpo,chen2024mallowspo,hong2024orpo}. DPO is a representative approach in this regard; by leveraging both preferred and non-preferred response, it enables the transfer of preference knowledge without relying on an explicit reward function.
RLHF has become a leading approach for aligning LLMs with human values, typically involving three stages: supervised fine-tuning, reward model training, and policy optimization using algorithms like PPO, GRPO, and REINFORCE++ \cite{ouyang2022training, casper2023open, achiam2023gpt}. Despite its effectiveness, RLHF’s reliance on multiple models and LLM sampling increases complexity \cite{zhao2023slic, casper2023open}. To address this, efficient offline alternatives like DPO have emerged, which transfer preference knowledge using both preferred and non-preferred responses without explicit reward models \cite{rafailov2023direct, xu2024contrastive, meng2024simpo}.

% However, recent studies have shown that DPO optimizes large language models (LLMs) based on preferences over entire responses, overlooking the varying importance of different tokens, which in turn limits its performance to some extent \cite{lin2024rho,zeng2024token,liu2024tis,gu2025mask}. Building on this line of research, our work investigates why the differences in token importance are critical to DPO training. We observe that many semantically similar tokens exist within the pairs used for DPO training. We hypothesize that training with semantically similar tokens in both positive and negative examples may lead to confusion—namely, the policy is required to output tokens with similar semantics while simultaneously being pushed to distinguish them. Based on this observation, we propose a simple yet effective solution, AAO, which enables the model itself to automatically re-weight tokens, thereby alleviating the confusion problem and significantly improving DPO performance. Furthermore, through both mathematical analysis and empirical experiments, we demonstrate the limiting effect of token confusion on DPO and reveal a novel phenomenon: semantic confusion may be one of the underlying causes of the "squeeze effect" in \cite{ren2024learning}.

Recent studies indicate DPO optimizes LLMs by considering whole-response preferences, ignoring token-level importance, which constrains its performance \cite{zeng2024token,liu2024tis,gu2025mask}. To address this, our work explores the impact of token importance in DPO training, observing that semantically similar tokens in positive and negative examples may cause confusion. 
To solve this, we introduce AAO, an approach that allows the model to re-weight tokens automatically, mitigating confusion and enhancing DPO performance. Our analysis further uncovers that semantic confusion may contribute to the “squeeze effect” identified in \cite{ren2024learning}.
%%%%%%%%%%%%%%%%%%%%%%%%%%%%%%%%%%%%%%%%%%%%%%%%%%%%%%%%%%%%%%%%%%%%%%%%%%%%%%%%%%%%%%%%%%%%%%%%%%%%%%%%%%%%%%%%%%%%%%%%
